\documentclass[../../../include/open-logic-section]{subfiles}

\begin{document}

\olfileid{sth}{ord-arithmetic}{mult}

\olsection{ક્રમવાચક સંખ્યાઓનો ગુણાકાર}

હવે ક્રમવાચક સંખ્યાઓના ગુણાકાર તરફ વળીએ. તેને પણ
સરવાળાની જેમ સમજીએ. ધારો કે $\alpha$ને~$\beta$થી
ગુણવું છે. પહોળાઈ $\alpha$ અને ઊંચાઈ $\beta$
હોય એવી લંબચોરસ જાળી કલ્પો: દરેક પંક્તિમાં આગળ
વધીએ, પછી બીજી પંક્તિમાં જઈએ\ldots એમ આખી
જાળીમાંથી પસાર થઈએ ત્યારે મળતો ક્રમપ્રકાર
$\alpha$ અને $\beta$નો ગુણાકાર છે. બીજી રીતે કહીએ
તો $\beta$ના \emph{દરેક} ઘટકની જગ્યાએ $\alpha$ની
એક નકલ મૂકવાથી $\alpha$ અને $\beta$નો ગુણાકાર
મળે છે.

તેને ઔપચારિક બનાવવા $\alpha$ અને $\beta$ના
કાર્ટેસિયન ગુણાકાર પર ઊલટો શબ્દકોશક્રમ વાપરીએ:

\begin{defn}
કોઈ પણ ક્રમવાચક સંખ્યાઓ $\alpha, \beta$ માટે
તેમનો ગુણાકાર $\alpha \ordtimes \beta = \ordtype{\alpha \times \beta, \rlexless}$ છે.
\end{defn}
\noindent
ફરી ખાતરી કરવી પડે કે આ વ્યાખ્યા સુગઠિત છે:

\begin{lem}\ollabel{ordtimeslessiswo}
કોઈ પણ ક્રમવાચક સંખ્યાઓ $\alpha$ અને $\beta$ માટે
$\tuple{\alpha \times \beta, \rlexless}$ સુક્રમ છે.
\end{lem}

\begin{proof}
બરાબર \olref[add]{ordsumlessiswo}ની જેમ.
\end{proof}
\noindent
ગુણાકાર નીચે પ્રમાણે વર્તે છે તે પણ સાબિત કરવું
અઘરું નથી:

\begin{lem}\ollabel{ordtimesrecursion}
કોઈ પણ ક્રમવાચક સંખ્યાઓ $\alpha, \beta$ માટે:
\begin{align*}
	\alpha \ordtimes 0 &= 0\\
	\alpha \ordtimes (\beta \ordplus 1) &=
		(\alpha \ordtimes \beta) \ordplus \alpha\\
	\alpha  \ordtimes \beta &=
		\supstrict_{\delta < \beta}(\alpha \ordtimes \delta) &&
		\text{જો $\beta$ સીમા ક્રમવાચક સંખ્યા હોય}.
\end{align*}
\end{lem}

\begin{proof}
સાબિતી સ્વાધ્યાય તરીકે છોડી છે.
\end{proof}

ખરેખર, સરવાળાની જેમ, સીધી વ્યાખ્યા આપવાને બદલે
આ પુનરાવર્તી સમીકરણોથી પણ ક્રમવાચક ગુણાકાર
વ્યાખ્યાયિત કરી શક્યા હોત. અને સરવાળાની જેમ
તેના કેટલાક ગુણધર્મો પણ પરિચિત છે:

\begin{lem}\ollabel{ordinalmultiplicationisnice}
$\alpha, \beta, \gamma$ ક્રમવાચક સંખ્યાઓ હોય, તો:
\begin{enumerate}
	\item\ollabel{ordtimes1} જો $\alpha \neq 0$ અને $\beta < \gamma$
	હોય, તો $\alpha \ordtimes \beta < \alpha \ordtimes \gamma$;
	\item\ollabel{ordtimes2} જો $\alpha \neq 0$ અને $\alpha \ordtimes
	\beta = \alpha\ordtimes\gamma$ હોય, તો $\beta = \gamma$;
	\item\ollabel{ordtimes3} $\alpha \ordtimes (\beta \ordtimes
	\gamma) = (\alpha \ordtimes \beta) \ordtimes \gamma$;
	\item\ollabel{ordtimes4} જો $\alpha \leq \beta$ હોય, તો $\alpha
	\ordtimes \gamma \leq \beta \ordtimes \gamma$;
	\item\ollabel{ordtimes5} $\alpha \ordtimes (\beta \ordplus
	\gamma) = (\alpha \ordtimes \beta )\ordplus (\alpha\ordtimes
	\gamma)$.
\end{enumerate}
\end{lem}

\begin{proof}
સાબિતી સ્વાધ્યાય તરીકે છોડી છે.
\end{proof}

ઇચ્છા મુજબ બીજા પરિણામો પણ સાબિત કરી શકો (અથવા
શોધીને વાંચી શકો). પરંતુ
\olref[ord-arithmetic][add]{ordsumnotcommute}ને
જોતાં નીચેનું આશ્ચર્યજનક ન લાગવું જોઈએ:

\begin{prop}
ક્રમવાચક ગુણાકાર અદલબદલીય નથી:
$2 \ordtimes \omega  =
\omega < \omega \ordtimes 2$
\end{prop}

\begin{proof}
$2 \ordtimes \omega = \supstrict_{n < \omega}(2\ordtimes  n) = \omega \in \supstrict_{n < \omega}(\omega \ordplus n) = \omega \ordplus \omega = \omega \ordtimes 2$.
\end{proof}
\noindent
ફરી, તેની સહજ સમજ સીધી છે. $2
\ordtimes \omega$ ગણવા દરેક પ્રાકૃતિક સંખ્યાની
જગ્યાએ બે ઘટકો મૂકો. એ જ ક્રમપ્રકાર પ્રાકૃતિક
સંખ્યાઓ વચ્ચે બધી ``અડધી'' સંખ્યાઓ ઉમેરીને,
એટલે $\Setabs{\nicefrac{n}{2}}{n \in \omega}$
પરનો સ્વાભાવિક ક્રમ વિચારીને, પણ મળે છે. આમ
જોતાં તેનો ક્રમપ્રકાર $\omega$ જેવો જ છે તે સ્પષ્ટ
થાય છે. પરંતુ $\omega \ordtimes 2$ ગણવા
$\omega$ની બે નકલો એક પછી એક મૂકો.

\begin{prob}
\olref[sth][ord-arithmetic][mult]{ordtimeslessiswo},
\olref[sth][ord-arithmetic][mult]{ordtimesrecursion}
અને
\olref[sth][ord-arithmetic][mult]{ordinalmultiplicationisnice}
સાબિત કરો.
\end{prob}

\end{document}
